\chapter{Talete di Mileto (Θαλῆς ὁ Μιλήσιος)}

\section{Le fonti}

\subsection{Le testimonianze antiche}

Talete di Mileto è tradizionalmente considerato il primo filosofo della cultura greca. Questa posizione non deriva da una testimonianza diretta della sua attività filosofica, ma soprattutto dalla ricostruzione elaborata da Aristotele.\footnote{Sull'interpretazione aristotelica del primato di Talete, cfr.\ \cite[17--19]{palumbo2015}.}

Le nostre conoscenze dipendono interamente dalla tradizione indiretta. Le testimonianze più antiche conservate risalgono a Erodoto, nel V secolo a.C., e a Platone e Aristotele, nel IV. Altre informazioni provengono da autori successivi, tra i quali Seneca, Diogene Laerzio e Proclo.

Questi autori tramandano notizie di natura differente. Erodoto riferisce episodi relativi alla vita pubblica e alle conoscenze scientifiche di Talete. Platone ne utilizza la figura per rappresentare un particolare atteggiamento del filosofo. Aristotele gli attribuisce alcune dottrine e le inserisce nella propria ricostruzione della storia del pensiero.

La distinzione tra queste testimonianze è essenziale per ricostruire la figura di Talete senza confondere le notizie tramandate con le interpretazioni degli autori che le hanno conservate.

\subsection{La documentazione delle dottrine}

Le testimonianze su Talete sono raccolte nel capitolo 11 dei \emph{Fragmente der Vorsokratiker} di Hermann Diels e Walter Kranz. La lettera A identifica i documenti che tramandano informazioni sulla sua vita e sul suo pensiero.

La testimonianza fondamentale è il passo della \emph{Metafisica} di Aristotele (I, 3, 983b6--27 = DK 11 A12), nel quale Talete viene presentato come iniziatore dell'indagine sui principi e gli viene attribuita la dottrina dell'acqua.

Non possediamo un testo autentico che permetta di stabilire come Talete formulasse e argomentasse questa dottrina. La ricostruzione del suo pensiero deve pertanto distinguere ciò che le fonti gli attribuiscono dalle categorie filosofiche utilizzate per interpretarlo.

\section{La vita}

\subsection{Profilo biografico}

Talete visse a Mileto, città greca della Ionia, tra il VII e il VI secolo a.C. La data di nascita tradizionalmente collocata intorno al 625 a.C. non è un dato documentato, ma il risultato di una ricostruzione cronologica. Anche la data della morte è incerta. Le notizie biografiche più articolate provengono da autori vissuti molti secoli dopo di lui, fra i quali Diogene Laerzio.

Gli antichi attribuirono a Talete interessi per la politica, la geometria e l'astronomia. Non possediamo opere autentiche che consentano di ricostruire direttamente la sua formazione e le sue attività; già le fonti antiche discordavano sull'esistenza di suoi scritti.

\begin{testimonianza}{Diogene Laerzio, \emph{Vite dei filosofi}, I, 37--38}{test:talete:cronologia}
Apollodoro, nella \emph{Cronologia}, afferma che nacque nel primo anno della trentacinquesima Olimpiade. Morì all'età di settantotto anni, oppure, secondo Sosicrate, di novanta; morì infatti nella cinquantottesima Olimpiade, all'epoca di Creso.
\end{testimonianza}

\noindent\textbf{Analisi.}
La testimonianza non contiene una data verificata sulla base di documenti contemporanei. Diogene Laerzio riunisce notizie che attribuisce a cronografi precedenti. L'indicazione della trentacinquesima Olimpiade corrisponde al 640 a.C.; nella tradizione testuale compare anche la variante della trentanovesima, corrispondente al 624 a.C. La cinquantottesima Olimpiade abbraccia invece gli anni 548--545 a.C. La differenza fra le due cronologie e fra le età di morte riferite da Diogene impedisce di assumere come certi gli anni riportati nelle ricostruzioni moderne.

L'\emph{akmḗ}, cioè l'età della piena maturità, veniva convenzionalmente collocata dai cronografi intorno ai quarant'anni. Il riferimento all'eclissi del 585 a.C. può quindi servire a costruire un'ipotesi cronologica, non a stabilire una data biografica indipendente. La conclusione storicamente sostenibile è più modesta: Talete appartiene alla prima metà del VI secolo a.C.; gli estremi della sua vita restano approssimativi.

\subsection{Le testimonianze di Erodoto}

Erodoto, vissuto nel V secolo a.C., costituisce la nostra fonte narrativa più antica per alcune attività pubbliche attribuite a Talete. Ne parla in relazione alla situazione politica della Ionia, a un'eclissi verificatasi durante una guerra e all'attraversamento del fiume Halys. Queste notizie richiedono però valutazioni distinte.

\noindent\textbf{La proposta politica.}
Secondo Erodoto, Talete suggerì alle città ioniche di istituire un consiglio comune con sede a Teo, conservando gli insediamenti delle singole città ma modificandone l'autonomia politica.

\begin{testimonianza}{Erodoto, \emph{Storie}, I, 170}{test:talete:teo}
Prima che la Ionia fosse sottomessa, Talete di Mileto, di discendenza fenicia, diede agli Ioni un consiglio utile: istituire un unico consiglio con sede a Teo, poiché Teo si trovava al centro della Ionia; le altre città avrebbero continuato a essere abitate, ma sarebbero state considerate come demi.
\end{testimonianza}

\noindent\textbf{Analisi.}
L'espressione greca \emph{hèn bouleutḗrion} (\textgreek{ἓν βουλευτήριον}) indica un unico organismo deliberativo. La proposta non consiste dunque in una generica alleanza: Erodoto le attribuisce una forma istituzionale definita. Il paragone con i demi indica che le città avrebbero conservato la propria esistenza come comunità locali, rinunciando però a una parte essenziale della loro indipendenza politica.

L'affermazione relativa alla discendenza fenicia di Talete è una notizia riferita dallo storico, non un fatto verificabile indipendentemente. Analogamente, il passo attesta che Erodoto conosceva una tradizione sulla proposta politica di Talete, ma non che tale proposta sia stata effettivamente discussa o realizzata.

\noindent\textbf{L'eclissi.}
Durante uno scontro tra Lidi e Medi, racconta Erodoto, il giorno si trasformò improvvisamente in notte. Lo storico attribuisce a Talete una precedente previsione del fenomeno.

\begin{testimonianza}{Erodoto, \emph{Storie}, I, 74}{test:talete:eclissi}
Mentre combattevano, accadde che il giorno diventasse improvvisamente notte. Talete di Mileto aveva preannunciato agli Ioni questo mutamento del giorno, indicando come limite proprio l'anno nel quale esso avvenne.
\end{testimonianza}

\noindent\textbf{Analisi.}
L'espressione greca \emph{oûron prothémenos eniautón} (\textgreek{οὖρον προθέμενος ἐνιαυτόν}) indica un limite temporale fissato nell'anno. Erodoto non afferma che Talete avesse previsto il giorno e l'ora dell'eclissi, né spiega il procedimento seguito.

L'episodio viene generalmente identificato con l'eclissi del 28 maggio 585 a.C.; tuttavia l'identificazione astronomica del fenomeno e l'attendibilità della previsione attribuita a Talete costituiscono due problemi diversi. La testimonianza documenta soprattutto la fama astronomica del milesio nella tradizione conosciuta da Erodoto. Non permette di dedurre che disponesse di un metodo esatto per prevedere le eclissi solari.

\noindent\textbf{Il fiume Halys.}
Una terza notizia riguarda il passaggio dell'esercito di Creso. Erodoto riferisce una tradizione che attribuiva a Talete la deviazione del fiume, ma esprime una diversa opinione su come l'attraversamento sarebbe avvenuto.

\begin{testimonianza}{Erodoto, \emph{Storie}, I, 75}{test:talete:halys}
Quando Creso giunse al fiume Halys, fece attraversare l'esercito, a mio parere, sui ponti che allora vi si trovavano; secondo il racconto comunemente diffuso fra i Greci, invece, fu Talete di Mileto a farlo passare.
\end{testimonianza}

\noindent\textbf{Analisi.}
Erodoto distingue espressamente la propria ricostruzione dal racconto circolante fra i Greci. Il seguito del passo descrive l'espediente attribuito a Talete: scavare un canale che deviasse parte del corso del fiume e ne rendesse possibile il guado.

Lo storico conserva dunque la notizia di una tradizione tecnica, senza condividerne necessariamente la veridicità. In termini documentari, possiamo affermare che a Talete veniva attribuita la capacità di risolvere un problema idraulico; non possiamo presentare la deviazione dell'Halys come un evento accertato. Questa distinzione fra notizia tramandata e giudizio del testimone è essenziale per utilizzare correttamente una fonte storica.

\subsection{Talete nella tradizione filosofica}

La figura di Talete assume un significato ulteriore nei testi di Platone e Aristotele. In questi autori gli episodi della sua vita non sono riferiti principalmente per costruirne una biografia, ma diventano esempi attraverso i quali discutere il carattere della sapienza e il valore della filosofia.

\noindent\textbf{I Sette Sapienti.}
Platone inserisce Talete fra i Sette Sapienti, associati a una sapienza espressa anche attraverso massime brevi e memorabili.

\begin{testimonianza}{Platone, \emph{Protagora}, 343a}{test:talete:sapienti}
Fra costoro vi erano Talete di Mileto, Pittaco di Mitilene, Biante di Priene, Solone della nostra città, Cleobulo di Lindo, Misone di Chene e, come settimo, Chilone di Sparta. Tutti erano ammiratori, seguaci e discepoli dell'educazione spartana: la loro sapienza si riconosceva nelle parole brevi e degne di essere ricordate pronunciate da ciascuno.
\end{testimonianza}

\noindent\textbf{Analisi.}
L'elenco colloca Talete in una tradizione culturale più ampia della filosofia intesa come indagine teorica. La sapienza dei sette viene presentata attraverso la brevità degli enunciati e il valore attribuito all'educazione spartana. Il passo documenta il ruolo attribuito a Talete nella memoria greca, non una sua dottrina particolare.

Anche le liste dei Sapienti non sono uniformi: la presenza di Misone nell'elenco platonico ricorda che non siamo davanti a un gruppo storicamente costituito secondo criteri immutabili. Il riconoscimento di Talete come sapiente va dunque distinto dal primato che gli attribuirà Aristotele nella ricerca dei principi naturali.

\noindent\textbf{La caduta nel pozzo.}
In un altro dialogo Platone contrappone la ricerca filosofica alle occupazioni immediate della vita quotidiana.

\begin{testimonianza}{Platone, \emph{Teeteto}, 174a}{test:talete:pozzo}
Talete, mentre osservava gli astri e guardava verso l'alto, cadde in un pozzo. Una serva tracia, arguta e spiritosa, lo derise perché desiderava conoscere le cose del cielo, ma gli sfuggivano quelle che si trovavano davanti a lui, ai suoi piedi.
\end{testimonianza}

\noindent\textbf{Analisi.}
Il passo appartiene al discorso di Socrate sulla distanza fra la vita filosofica e quella di chi si occupa soltanto degli affari immediati. L'aneddoto è seguito dalla sua estensione alla figura del filosofo in generale. Questa generalizzazione ne chiarisce la funzione: Talete rappresenta l'uomo che rivolge l'attenzione a ciò che oltrepassa l'esperienza quotidiana e per questo può apparire ridicolo nelle occupazioni più semplici.

Il testo non è quindi una prova indipendente che Talete sia realmente caduto in un pozzo. È invece una fonte decisiva per comprendere come Platone costruisca l'immagine del filosofo. Il carattere interpretativo della testimonianza platonica è una delle difficoltà metodologiche sottolineate da Palumbo.\footnote{Sulla mediazione interpretativa delle testimonianze platoniche, cfr.\ \cite[17]{palumbo2015}.}

\noindent\textbf{L'affare dei frantoi.}
Aristotele utilizza un episodio di segno diverso: Talete avrebbe sfruttato le proprie conoscenze astronomiche per anticipare un abbondante raccolto di olive e assicurarsi, a basso prezzo, l'uso dei frantoi.

\begin{testimonianza}{Aristotele, \emph{Politica}, I, 11, 1259a6--20 (DK 11 A10)}{test:talete:frantoi}
Raccontano che, essendo rimproverato per la sua povertà, quasi che la filosofia fosse inutile, Talete avesse previsto grazie all'astronomia, quando era ancora inverno, un abbondante raccolto di olive. Con il poco denaro disponibile versò anticipi per tutti i frantoi di Mileto e di Chio, ottenendoli a buon prezzo perché nessuno gareggiava con lui. Quando giunse il momento del raccolto e molti li richiesero contemporaneamente, li concesse in uso al prezzo che volle e raccolse molto denaro. Mostrò così che per i filosofi è facile arricchirsi, se lo desiderano, ma che non è questo il loro interesse.
\end{testimonianza}

\noindent\textbf{Analisi.}
Il racconto è inserito nella trattazione aristotelica dell'acquisizione della ricchezza. Il sapere astronomico fornisce la premessa dell'operazione: Talete avrebbe previsto una futura crescita della domanda. Il guadagno deriverebbe però dall'aver acquistato in anticipo la disponibilità dei frantoi, non dalla previsione in sé. Aristotele presenta l'espediente come un caso di monopolio, riconoscendovi un principio economico applicabile anche in altre circostanze.

La funzione filosofica dell'episodio risulta dalla sua conclusione. Alla critica secondo cui la filosofia sarebbe inutile perché non procura ricchezza si risponde che la capacità di arricchirsi non coincide con il fine del filosofo. Talete può utilizzare il proprio sapere per conseguire un vantaggio materiale, ma non pone quel vantaggio al centro della sua attività.

Il testo è introdotto come una storia riferita da altri, non come un resoconto biografico verificato. Accostato al racconto platonico del pozzo (testimonianza~\ref{test:talete:pozzo}), rende visibili due rappresentazioni complementari del filosofo: apparentemente inadatto alla vita ordinaria, eppure capace di servirsi efficacemente delle proprie conoscenze. La testimonianza aristotelica è riportata anche da Palumbo nella discussione sul primato filosofico di Talete.\footnote{Per la presentazione della testimonianza di Aristotele, \cite[18]{palumbo2015}.}

\subsection{Le tradizioni successive}

Le testimonianze più tarde attribuiscono a Talete anche l'introduzione della geometria egiziana in Grecia e la capacità di risolvere problemi di misura. Queste notizie documentano la formazione di una tradizione sul sapere matematico del milesio; la distanza cronologica impedisce di considerarle automaticamente resoconti attendibili dei singoli episodi.

\noindent\textbf{Il viaggio in Egitto.}
Proclo colloca Talete all'inizio della storia della geometria greca, in una ricostruzione che viene tradizionalmente collegata anche ai materiali storici di Eudemo di Rodi.

\begin{testimonianza}{Proclo, \emph{Commento al primo libro degli Elementi di Euclide}, p.~65, ed.\ Friedlein}{test:talete:egitto}
Talete, recatosi in Egitto, portò per primo in Grecia questa conoscenza. Fece egli stesso numerose scoperte e trasmise ai successori i principi di molte altre, affrontando alcuni problemi in modo più generale e altri in modo più empirico.
\end{testimonianza}

\noindent\textbf{Analisi.}
Proclo non sta scrivendo una biografia di Talete, ma ripercorre lo sviluppo storico di una disciplina. Collega l'origine egiziana della geometria alle esigenze pratiche della misurazione dei terreni e attribuisce a Talete il trasferimento di queste conoscenze nel mondo greco.

La distinzione fra procedimenti più generali e procedimenti empirici è particolarmente interessante: nel resoconto di Proclo, Talete si colloca fra l'impiego pratico delle conoscenze geometriche e la loro elaborazione teorica. Ciò non dimostra tuttavia che egli disponesse già delle forme di dimostrazione proprie della geometria successiva. La testimonianza documenta soprattutto il ruolo che la tradizione storiografica gli attribuiva nella storia della matematica.

\noindent\textbf{La misura delle piramidi.}
Diogene Laerzio riporta una notizia attribuita a Geronimo di Rodi.

\begin{testimonianza}{Diogene Laerzio, \emph{Vite dei filosofi}, I, 27}{test:talete:piramidi}
Geronimo riferisce che Talete misurò l'altezza delle piramidi mediante la loro ombra, osservando il momento in cui la nostra ombra è lunga quanto noi stessi.
\end{testimonianza}

\noindent\textbf{Analisi.}
La procedura descritta presuppone una relazione geometrica elementare: in un dato momento, quando la lunghezza dell'ombra di un corpo verticale coincide con la sua altezza, lo stesso rapporto vale per altri corpi verticali illuminati nelle medesime condizioni. L'altezza della piramide potrebbe così essere determinata misurandone l'ombra, tenendo conto delle condizioni geometriche effettive della misura.

Non è tuttavia possibile passare da questa descrizione all'affermazione che Talete avesse formulato e dimostrato un teorema generale sulla similitudine dei triangoli. L'applicazione pratica di una regolarità e la sua dimostrazione astratta sono risultati differenti.

Come nella testimonianza di Proclo (testimonianza~\ref{test:talete:egitto}), la notizia è significativa per ricostruire l'immagine di Talete trasmessa dalla storia antica della geometria. La sua attendibilità biografica rimane distinta dal significato matematico del procedimento narrato.

\section{Le opere}

\subsection{Gli scritti attribuiti a Talete}

L'esistenza di opere scritte da Talete era già discussa nell'antichità. Diogene Laerzio (I, 23) riferisce tradizioni differenti: secondo alcune, Talete non avrebbe lasciato alcuno scritto; secondo altre, sarebbe stato autore di due trattati, \emph{Sul solstizio} e \emph{Sull'equinozio}.

Al suo nome fu associata anche un'opera di astronomia nautica, attribuita però da alcune fonti a Foco di Samo.

La tradizione conserva inoltre lettere indirizzate ad altri sapienti, tra i quali Solone e Ferecide. Tali documenti sono generalmente considerati composizioni posteriori.

Nessuna delle opere attribuite a Talete è giunta fino a noi in forma autentica.

\subsection{I limiti della ricostruzione}

L'assenza di scritti impedisce di conoscere la forma originaria dell'esposizione di Talete. Non sappiamo quali termini utilizzasse per esprimere le proprie concezioni, come organizzasse gli argomenti o quali ragioni adducesse a sostegno delle dottrine che gli vengono attribuite.

Questa difficoltà riguarda soprattutto la tesi dell'acqua. La conosciamo attraverso Aristotele, ma non possediamo un testo di Talete che permetta di confrontare la testimonianza aristotelica con la formulazione originaria.

La sua filosofia deve pertanto essere ricostruita attraverso un numero limitato di testimonianze, senza presupporre l'esistenza di un sistema dottrinale compiuto.

\section{Il pensiero}

\subsection{L'indagine sulla natura}

La tradizione attribuisce a Talete un interesse per l'osservazione dei fenomeni naturali e per la ricerca delle cause che ne spiegano il comportamento.

Le testimonianze riguardano fenomeni differenti: la posizione della Terra, i terremoti, il movimento dei corpi, gli astri e le proprietà dell'acqua. In ciascuno di questi ambiti, Talete viene presentato come un pensatore interessato a comprendere i fenomeni attraverso relazioni appartenenti alla natura stessa.

Questo orientamento costituisce uno dei motivi del primato filosofico riconosciutogli da Aristotele.

Non è possibile, tuttavia, ricostruire con certezza un metodo di ricerca uniforme. Le attribuzioni sono frammentarie e provengono da autori che operarono in epoche successive. Esse documentano soprattutto l'immagine di Talete come indagatore della natura, senza consentire di definirne precisamente le procedure conoscitive.

\subsection{L'acqua come principio delle cose}

La testimonianza fondamentale sul pensiero di Talete è contenuta nella \emph{Metafisica} di Aristotele (I, 3, 983b6--27 = DK 11 A12).

Aristotele esamina le dottrine dei primi pensatori alla luce del problema delle cause e dei principi. Secondo la sua ricostruzione, i filosofi più antichi avrebbero individuato soprattutto principi materiali, cioè elementi dai quali le cose derivano e dei quali sono costituite.

A Talete attribuisce l'identificazione di questo principio con l'acqua.

Il termine \emph{archḗ} (ἀρχή) significa «inizio», «origine» o «principio». Nel quadro teorico aristotelico, il principio materiale è ciò di cui le cose sono costituite, ciò da cui provengono e ciò in cui si dissolvono.

Questa definizione presuppone che, attraverso il mutamento delle cose, permanga qualcosa che ne costituisce il fondamento. I singoli enti nascono, si trasformano e cessano di essere, mentre il principio materiale permette di comprenderne l'origine e la costituzione.

Aristotele interpreta la dottrina dell'acqua entro questa concezione. Non sappiamo, però, se Talete considerasse l'acqua un sostrato materiale permanente, nel significato che questa nozione assume nella metafisica aristotelica.

È possibile che l'acqua fosse intesa anzitutto come origine delle cose, senza che venisse elaborata una teoria generale della sostanza. La documentazione non consente di stabilirlo con certezza.

\subsection{Le osservazioni sull'umidità}

Aristotele propone alcune possibili ragioni per spiegare perché Talete avrebbe identificato nell'acqua il principio delle cose.

Il nutrimento degli esseri viventi è umido; il calore sembra svilupparsi dall'umidità e conservarsi grazie a essa; i semi dai quali nascono gli organismi possiedono una natura umida. Da queste osservazioni Talete avrebbe potuto ricavare l'importanza dell'acqua nei processi naturali.

Il ragionamento procede dai fenomeni particolari verso una spiegazione più generale. Proprietà osservate in situazioni differenti vengono ricondotte a un elemento comune.

La generalizzazione presenta tuttavia una difficoltà. Constatare che l'acqua interviene in numerosi processi naturali non equivale a dimostrare che essa costituisca il principio di tutte le cose.

Aristotele, inoltre, introduce queste ragioni in forma esplicitamente congetturale. Attribuisce a Talete la dottrina dell'acqua, ma non afferma di conoscere con certezza il procedimento attraverso il quale il filosofo vi sarebbe giunto.

La distinzione tra dottrina attribuita e spiegazione congetturale è decisiva: le osservazioni sull'umidità non sono argomenti conservati in uno scritto di Talete, ma ipotesi avanzate da Aristotele per comprenderne il pensiero.\footnote{Sul carattere congetturale delle spiegazioni aristoteliche, cfr.\ Aristotele, \emph{Metafisica}, I, 3, 983b20--27, e \cite[18--19]{palumbo2015}.}

\subsection{La terra sostenuta dall'acqua}

Un'altra dottrina attribuita a Talete riguarda la posizione e la stabilità della Terra.

Aristotele riferisce che, secondo un'antica concezione associata a Talete, la Terra poggia sull'acqua e vi rimane sostenuta come un corpo galleggiante (\emph{De caelo}, II, 13, 294a28--34).

Seneca tramanda una spiegazione collegata a questa concezione (\emph{Questioni naturali}, III, 14 = DK 11 A15). La Terra sarebbe sostenuta dall'acqua e, quando questa si muove, oscillerebbe come un'imbarcazione, provocando i terremoti.

Il ragionamento ricorre a un'analogia tra due fenomeni. Il movimento di un'imbarcazione sull'acqua costituisce un caso osservabile; il movimento della Terra viene interpretato mediante una relazione simile.

L'ipotesi riconduce così il terremoto a una causa fisica. Non è necessario attribuirlo all'azione di una divinità: esso può essere spiegato come conseguenza del movimento dell'acqua.

Questa testimonianza non coincide necessariamente con la dottrina aristotelica dell'acqua come principio universale.

Affermare che l'acqua sostiene la Terra significa attribuirle una funzione cosmologica particolare. Affermare che l'acqua è il principio di tutte le cose significa invece attribuirle una funzione generale di origine e fondamento della realtà.

La prima proposizione non implica logicamente la seconda. La distinzione impedisce di trasformare una spiegazione fisica circoscritta in una dottrina metafisica sistematica senza ulteriori testimonianze.

\subsection{L'anima, il movimento e il divino}

Aristotele attribuisce a Talete alcune considerazioni riguardanti l'anima e il movimento.

Nel \emph{De anima} (I, 2, 405a19--21) riferisce che Talete avrebbe riconosciuto un'anima persino nella pietra magnetica, poiché questa possiede la capacità di muovere il ferro.

La notizia suggerisce una relazione tra l'anima e la capacità di produrre movimento. Aristotele, tuttavia, la presenta con cautela, richiamando quanto era stato tramandato sul pensiero di Talete.

In un altro passo del \emph{De anima} (I, 5, 411a7--8) Aristotele collega a Talete l'affermazione secondo cui tutte le cose sarebbero piene di dèi.

Queste testimonianze mostrano che la ricerca di spiegazioni naturali non comporta necessariamente l'esclusione del divino. Una concezione della natura può riconoscere principi di movimento e di vita interni alle cose, senza separarli nettamente dalla dimensione divina.

Non sappiamo, tuttavia, se Talete abbia elaborato una dottrina organica dell'anima o della divinità. Le testimonianze aristoteliche non sono sufficienti per attribuirgli con certezza una teoria secondo cui tutte le cose possiederebbero un'anima.

\subsection{Geometria e astronomia}

Le fonti antiche attribuiscono a Talete conoscenze geometriche e astronomiche.

Proclo, che conserva notizie risalenti alla tradizione matematica antica e in particolare a Eudemo di Rodi, gli assegna diversi risultati geometrici. Tra questi figurano proprietà del cerchio, degli angoli opposti al vertice e dei triangoli.

Al suo nome è tradizionalmente associato anche il risultato secondo cui un triangolo inscritto in una semicirconferenza, con un lato coincidente con il diametro, è rettangolo. La tradizione non permette tuttavia di stabilire se Talete ne abbia fornito una dimostrazione nel senso successivamente sviluppato dalla geometria greca.

I racconti sulla misurazione delle piramidi e delle distanze delle navi dalla costa collegano le sue conoscenze geometriche a problemi pratici.

In astronomia gli vengono attribuite osservazioni relative agli astri, ai solstizi e agli equinozi. L'attribuzione più celebre resta quella della previsione dell'eclissi del 585 a.C., la cui attendibilità è assai dubbia.

Queste notizie documentano la fama di Talete come conoscitore della geometria e dell'astronomia, ma non consentono di ricostruire con precisione i suoi procedimenti teorici.

\section{L'eredità}

\subsection{La tradizione milesia}

La tradizione storiografica colloca Talete all'origine della cosiddetta scuola di Mileto, alla quale vengono associati Anassimandro e Anassimene.

La denominazione non deve necessariamente far pensare a un'istituzione organizzata o a una successione di maestri e discepoli storicamente documentata. Indica soprattutto una continuità di interessi e di problemi.

I tre pensatori sono accomunati dalla ricerca di spiegazioni dei fenomeni naturali e dall'interrogazione sull'origine e sulla costituzione del mondo.

Secondo la ricostruzione tradizionale, Anassimandro avrebbe individuato il principio nell'\emph{ápeiron}, l'indefinito o illimitato, mentre Anassimene lo avrebbe identificato nell'aria.

La successione delle dottrine documenta, nella ricostruzione antica, un problema comune affrontato attraverso soluzioni differenti. Essa non dimostra però che tutti e tre i pensatori possedessero già una nozione di \emph{archḗ} pienamente corrispondente a quella aristotelica.

\subsection{Palumbo: il problema della falsa evidenza}

Lidia Palumbo individua nel pensiero di Talete un problema storiografico fondamentale. La sua posizione di primo filosofo appare consolidata dalla tradizione, ma dipende da una determinata interpretazione di ciò che costituisce la filosofia.

Il punto decisivo riguarda la ricostruzione aristotelica. Aristotele considera la filosofia come ricerca delle cause e dei principi e interpreta le dottrine dei predecessori attraverso questa prospettiva. Talete viene così presentato come l'iniziatore dell'indagine sull'\emph{archḗ}, identificata con l'acqua.

Palumbo osserva che probabilmente Talete non pensò mai all'acqua nei termini di un principio universale, come viene formulato dalla filosofia aristotelica. La questione non riguarda necessariamente l'autenticità della dottrina dell'acqua, ma il significato attribuito a tale dottrina nella ricostruzione successiva.

Il rischio consiste nel considerare originaria una categoria elaborata più tardi. La sistematicità dell'interpretazione aristotelica rende coerenti le dottrine dei predecessori, ma può anche nasconderne le differenze e le specificità.

In questo senso Palumbo richiama il compito dello storico della filosofia: «sostituire un problema ad una falsa evidenza».\footnote{\cite[18]{palumbo2015}. La formulazione è citata letteralmente dal testo.}

La designazione di Talete come primo filosofo diventa così un problema da indagare, anziché una classificazione da accettare senza esame. La ricostruzione storica deve distinguere, per quanto le fonti consentono, le concezioni attribuibili al pensatore arcaico dalle categorie attraverso le quali sono state successivamente interpretate.

La riflessione di Palumbo non intende negare il primato tradizionale di Talete, ma chiarirne il fondamento e i limiti. L'importanza della sua figura resta riconoscibile anche quando si rinuncia ad attribuirgli una dottrina sistematica dei principi.

\subsection{Nietzsche: l'intuizione dell'unità}

In \emph{La filosofia nell'epoca tragica dei Greci}, Friedrich Nietzsche propone un'interpretazione del primato filosofico di Talete.\footnote{L'interpretazione nietzscheana è qui ripresa attraverso \cite[19]{palumbo2015}.}

Secondo Nietzsche, la semplice osservazione delle proprietà dell'acqua e la formulazione di una spiegazione dei fenomeni naturali non sarebbero sufficienti a distinguere la filosofia dalla ricerca scientifica.

Il carattere propriamente filosofico della dottrina risiederebbe nell'idea che la molteplicità delle cose possa essere ricondotta a un'unica realtà fondamentale.

Nietzsche esprime questo significato mediante la formula «Tutto è uno». Non si tratta di una citazione autentica di Talete, ma di una formulazione interpretativa attraverso la quale Nietzsche individua il contenuto filosofico della sua ricerca.

Il passaggio dalla molteplicità dei fenomeni alla loro unità richiede, secondo Nietzsche, due capacità: l'astrazione e la fantasia.

L'astrazione consente di oltrepassare la considerazione isolata delle cose. La fantasia permette di concepire un'unità che non è direttamente offerta dall'esperienza.

In questa prospettiva, Talete non è il primo filosofo perché avrebbe scoperto una particolare proprietà dell'acqua, ma perché avrebbe concepito la possibilità di comprendere l'intera realtà mediante un unico principio.

\subsection{Severino: l'identità del molteplice}

Emanuele Severino approfondisce il significato filosofico della dottrina di Talete ponendo al centro il problema del Tutto.\footnote{\cite[48--50]{severino2010}.}

La filosofia, secondo Severino, nasce dall'interesse per la totalità della realtà. Le cose appaiono differenti e molteplici, ma appartengono a un medesimo orizzonte. Per pensare il Tutto occorre individuare ciò che permette di ricondurre la molteplicità a un'unità.

La questione fondamentale non consiste dunque nello stabilire semplicemente quale elemento materiale venga prima degli altri, ma nel comprendere che cosa costituisca il fondamento unitario delle cose.

Severino interpreta l'acqua di Talete all'interno di questo problema. Se l'acqua fosse soltanto quella sensibile, distinguibile dalla terra, dall'aria e dagli altri elementi, non potrebbe costituire ciò che è comune a tutte le cose.

L'acqua assume perciò, nella sua interpretazione, il valore di un'espressione ancora inadeguata dell'unità del reale. Talete avrebbe concepito un fondamento universale, ma lo avrebbe indicato con il nome di un elemento particolare.

Questa difficoltà permette a Severino di interpretare anche il passaggio ad Anassimandro, che individua il principio non più in una realtà determinata come l'acqua, ma nell'illimitato.

Si tratta di una ricostruzione filosofica moderna, non di una dottrina direttamente documentata negli scritti di Talete.

\subsection{Considerazioni conclusive}

Talete occupa una posizione particolare nella storia della filosofia: la tradizione lo considera il suo iniziatore, ma non possediamo alcun testo autentico che permetta di ricostruirne direttamente il pensiero.

Le testimonianze gli attribuiscono interessi naturalistici, matematici e astronomici. Aristotele lo presenta come l'iniziatore della ricerca dei principi, attribuendogli la dottrina dell'acqua.

Il significato filosofico di questa dottrina è stato diversamente interpretato. Aristotele la riconduce alla ricerca delle cause materiali; Nietzsche vi individua l'intuizione dell'unità del reale; Severino la interpreta come un primo tentativo di pensare il fondamento unitario del Tutto. Palumbo richiama invece l'attenzione sulla distanza tra le concezioni attribuibili a Talete e le categorie attraverso le quali sono state successivamente comprese.

Il primato filosofico di Talete risulta così inseparabile dalle testimonianze che ne hanno conservato la memoria e dalle interpretazioni che ne hanno determinato la fortuna.

% Il repertorio riguarda esclusivamente gli studi effettivamente citati
% nelle note di questo capitolo; le pagine puntuali restano nelle note.
\addsec{Bibliografia}
\printbibliography[heading=none]
